\begin{tikzpicture}[figure picture,x=1mm,y=1mm]
 \path[use as bounding box] (0,0) rectangle (169,77);
 \begin{scope}[overlay]
  \node[anchor=west] at (11,73) {轨迹：共享一个时间演化规律};
  \node[anchor=west] at (112,73) {截面：固定时刻的分布};
  \node[anchor=west] at (11,66) {$X_t$};
  \node[anchor=west] at (115,66) {$p_t(x)$};
  \begin{axis}[
   at={(11mm,19mm)},anchor=south west,width=87mm,height=44mm,
   scale only axis,xmin=0,xmax=4,ymin=-1.1,ymax=1.4,
   xtick={0,.5,2,4},ytick={-1,0,1},axis lines=left,grid=major,
   xlabel={时刻 $t$},
   xlabel style={at={(axis description cs:1,-.14)},anchor=north east}
  ]
   \addplot[FigureBlue!55!white,line width=.7pt]
    table[x=t,y=path2]{figures/math-preparation/reading-restructure/ch13/ou.tsv};
   \addplot[FigureBlue!55!white,line width=.7pt]
    table[x=t,y=path3]{figures/math-preparation/reading-restructure/ch13/ou.tsv};
   \addplot[FigureBlue!55!white,line width=.7pt]
    table[x=t,y=path4]{figures/math-preparation/reading-restructure/ch13/ou.tsv};
   \addplot[FigureBlue,line width=1pt]
    table[x=t,y=path1]{figures/math-preparation/reading-restructure/ch13/ou.tsv};
   \addplot[FigureInk,dashed,line width=.9pt]
    table[x=t,y=mean]{figures/math-preparation/reading-restructure/ch13/ou.tsv};
   \draw[FigureRed,densely dotted,line width=.7pt] (axis cs:.5,-1.1)--(axis cs:.5,1.4);
   \draw[FigureYellow,densely dotted,line width=.7pt] (axis cs:2,-1.1)--(axis cs:2,1.4);
  \end{axis}
  \begin{axis}[
   at={(115mm,19mm)},anchor=south west,width=47mm,height=44mm,
   scale only axis,xmin=-1.1,xmax=1.7,ymin=0,ymax=1.3,
   xtick={-1,0,1},ytick={0,.5,1},axis lines=left,grid=major,
   xlabel={状态 $x$},
   xlabel style={at={(axis description cs:1,-.14)},anchor=north east}
  ]
   \addplot[FigureRed,line width=1pt,domain=-1.1:1.7,samples=120]
    {exp(-(x-exp(-.5))^2/(.36*(1-exp(-1))))/sqrt(pi*.36*(1-exp(-1)))};
   \addplot[FigureYellow,line width=1pt,dashed,domain=-1.1:1.7,samples=120]
    {exp(-(x-exp(-2))^2/(.36*(1-exp(-4))))/sqrt(pi*.36*(1-exp(-4)))};
  \end{axis}
  \draw[FigureInk,dashed,line width=.9pt] (11,6)--(18,6);
  \node[anchor=west] at (19,6) {理论均值 $\mathrm e^{-t}$；四条精确网格模拟};
  \draw[FigureRed,line width=1pt] (112,6)--(118,6);
  \node[anchor=west] at (119,6) {$t=0.5$};
  \draw[FigureYellow,dashed,line width=1pt] (141,6)--(147,6);
  \node[anchor=west] at (148,6) {$t=2$};
 \end{scope}
\end{tikzpicture}
